\documentclass[10pt,a4paper]{amsart}
\usepackage[margin=19mm]{geometry}
\usepackage{amsmath,amssymb,mathtools}
\usepackage[scr=rsfs,cal=euler]{mathalfa}
\usepackage{xfrac}
\usepackage[unicode,hidelinks,bookmarks=false]{hyperref}
\newcommand{\ie}{i.e.\ }
\newcommand{\cf}{cf.\ }
\newcommand{\resp}{respectively\ }
\newcommand{\Subsection}{\subsection}
\providecommand{\nobreakdash}{\nobreak\hbox{-}}
\setlength{\emergencystretch}{2em}
\multlinegap0pt
\usepackage{bm}
\usepackage{bbm}
\usepackage{dsfont}
\usepackage{xspace}
\usepackage{cancel}
\usepackage{mathtools}
\usepackage{amsthm}
\makeatletter
\@ifundefined{theorem}{\newtheorem{theorem}{Theorem}}{}
\@ifundefined{lemma}{\newtheorem{lemma}[theorem]{Lemma}}{}
\makeatother
\title[A Convergent Geometry-Aware Reduction for Diffusion in Branc]{A Convergent Geometry-Aware Reduction for Diffusion in Branched Tubular Networks}
\author{Zachary M. Miksis, Gillian Queisser}
\date{Source: arXiv:2501.08247v4 (CC BY 4.0)}
\begin{document}
\maketitle

\noindent\emph{Source and licence.} A Convergent Geometry-Aware Reduction for Diffusion in Branched Tubular Networks. Version arXiv:2501.08247v4 (CC BY 4.0), distributed under the Creative Commons Attribution 4.0 International licence, as recorded in the arXiv licence metadata for this exact version and shown on the abstract page https://arxiv.org/abs/2501.08247v4. Full rights notice: licenses/arxiv-2501.08247-CC-BY-4.0.txt.\par\smallskip

\noindent\textit{Editorial scope.} Counted unit: the Lemma whose source label is \texttt{stab-lemma} in the article (source0.tex lines 317--327, source label \texttt{stab-lemma}), delivered together with the article's own complete original proof of it (source0.tex lines 329--409, one \texttt{proof} environment of 81 lines). No mathematics is written, repaired or re-derived: statement and proof are the article's own text line for line. Required context retained uncounted: eq:dFJ (lines 245--247). Every definition, display and label of those lines is retained unchanged. Omitted from this package and not redistributed: 1--244, 248--316, 328--328, 410--763. Nothing inside a retained excerpt is omitted or altered: every retained body is delivered exactly as the article's lines are written, so no omission receipt of any kind accompanies this package. Citations: the retained excerpts carry no citation command at all, so no bibliography entry of the article is transcribed and no reference is redistributed. Reference adaptation: none was needed and none was made; every reference inside the delivered unit resolves inside it, so no unresolved reference or citation is produced. Printed slips kept literal: no printed word, symbol, hypothesis or number of the counted statement or of its proof was altered. Layout and class: the article's own class is \texttt{article} with packages including graphicx, amsmath, amssymb, amsthm, hyperref, geometry, xcolor, caption, subcaption, lineno, natbib, multirow; the wrapper is the lane's pinned \texttt{amsart} editorial wrapper at 10pt on A4, which restates the theorem-like environments the retained text needs and adds the article's own macro definitions that the retained text needs (none), with extra wrapper packages (bm, bbm, dsfont, xspace, cancel, mathtools). The delivered numbering is the unit's own. Assets: no retained span contains a figure environment, a graphics-inclusion command, a PDF-page inclusion, a table or a TikZ picture; no image file is copied into this package, the package declares no asset dependency and no image file is redistributed with it. Licence: version arXiv:2501.08247v4 of this article is distributed under the Creative Commons Attribution 4.0 International licence, as recorded in the arXiv licence metadata for this exact version and shown on the abstract page https://arxiv.org/abs/2501.08247v4. A full rights notice is delivered at licenses/arxiv-2501.08247-CC-BY-4.0.txt. Characters: every retained excerpt is pure ASCII, so the delivered engine, XeLaTeX with its default Latin Modern text font, reports no missing character.
\begin{equation}
    \mathcal{A}(x)\frac{\partial u}{\partial t} = w(x) \frac{D_0}{R^2} \frac{\partial}{\partial x} \mathcal{F}(u) + [1 - w(x)] D_0 \frac{\partial^2 u}{\partial x^2} + \frac{\mathcal{S}}{\pi\Delta x R^2} \mathcal{J} \label{eq:dFJ}
\end{equation}

\begin{lemma}[Discrete energy stability for the expansion model]\label{stab-lemma}
    Consider the expansion-model flux $\mathcal{F}(u) = a(x) \partial_x u + b(x) \partial_x^3 u$ on $[0, L]$ with
    \begin{equation}
        a(x) \;=\; R(x)^2 + \frac{\Delta x^2}{4} R_x(x)^2, \qquad b(x) \;=\; \frac{\Delta x^2}{8} R(x)^2, \label{eq:expansion-ab}
    \end{equation}
    and the finite-volume semi-discretization of $\partial_t u = (D_0/R^2)\, \partial_x \mathcal{F}(u)$ on a uniform grid of spacing $\Delta x$, with no-flux boundary conditions $F_{-1/2} = F_{n+1/2} = 0$ and numerical conditions $(u_0)_{xx} = (u_n)_{xx} = 0$. Suppose
    \begin{equation}
        \Delta x \;\leq\; \Delta x_*, \qquad \Delta x_* \;:=\; \frac{R_{\min}^2}{\|R\|_\infty^2 + \|R\|_\infty \|R_x\|_\infty}. \label{eq:stab-threshold}
    \end{equation}
    Given geometry (A1), the discretization is discretely $L^2$ energy stable in the area-weighted norm $\|u\|_R^2 := \sum_i u_i^2 R_i^2 \Delta x$: $dE/dt \leq 0$ for the homogeneous problem, where $E := \tfrac{1}{2} \|u\|_R^2$.
\end{lemma}

\begin{proof}
   The proof has three parts: setup with the area-weighted energy; summation by parts and decomposition of the third-derivative flux contribution; and a combined estimate that closes via the threshold \eqref{eq:stab-threshold}.

    \emph{Setup.}
    The full equation \eqref{eq:dFJ} with $w(x) \equiv 1$ reads $\mathcal{A}(x)\, \partial_t u = (D_0/R^2)\, \partial_x \mathcal{F}(u)$, where $\mathcal{A}(x) = 1 + (\Delta x^2/12)(R_x/R)^2$ is bounded with $1 \le \mathcal{A}(x) \le 1 + (\Delta x^2/12)(\|R_x\|_\infty/R_{\min})^2$, in particular uniformly bounded above and below. We absorb $\mathcal{A}$ into the time variable by a positive rescaling that does not affect the sign of $dE/dt$, setting $\mathcal{A} \equiv 1$ for the stability analysis without loss of generality. The semi-discrete scheme is
    \begin{equation}
        \dot{u}_i = \frac{D_0}{R_i^2}\, \frac{F_{i+1/2} - F_{i-1/2}}{\Delta x},
    \end{equation}
    with $F_{-1/2} = F_{n+1/2} = 0$. Using the area-weighted inner product $\langle u, v\rangle_R := \sum_i u_i v_i R_i^2 \Delta x$ and the corresponding energy $E := \tfrac{1}{2}\langle u, u\rangle_R$, the $R^2$ weight cancels the $D_0/R_i^2$ factor on the right-hand side:
    \begin{equation}
        \frac{dE}{dt} \;=\; \sum_i u_i\, R_i^2 \, \dot u_i\, \Delta x \;=\; D_0 \sum_i u_i (F_{i+1/2} - F_{i-1/2}).
    \end{equation}
    Using $F_{-1/2} = F_{n+1/2} = 0$ to discard boundary terms, summation by parts gives
    \begin{equation}
        \frac{dE}{dt} \;=\; D_0 \sum_i (u_i - u_{i+1})\, F_{i+1/2} \;=\; -D_0 \sum_i \Delta x\, (u_{i+1/2})_x\, F_{i+1/2},
    \end{equation}
    where $(u_{i+1/2})_x := (u_{i+1} - u_i)/\Delta x$. Substituting $F_{i+1/2} = a_{i+1/2}(u_{i+1/2})_x + b_{i+1/2}(u_{i+1/2})_{xxx}$,
    \begin{equation}
        \frac{dE}{dt} \;=\; \underbrace{-D_0 \Delta x \sum_i a_{i+1/2}\, (u_{i+1/2})_x^2}_{T_1} \;+\; \underbrace{-D_0 \Delta x \sum_i b_{i+1/2}\, (u_{i+1/2})_x\, (u_{i+1/2})_{xxx}}_{T_2}.
    \end{equation}
    The first sum is non-positive, $T_1 \leq 0$. The second sum has indefinite sign in general.

    \emph{Decomposition of $T_2$.}
    Set $D^+ u_i := (u_{i+1} - u_i)/\Delta x$ and $(u_i)_{xx} := (u_{i+1} - 2u_i + u_{i-1})/\Delta x^2$. By definition of the third-derivative finite difference,
    \begin{equation}
        (u_{i+1/2})_{xxx} \;=\; \frac{(u_{i+1})_{xx} - (u_i)_{xx}}{\Delta x}.
    \end{equation}
    Substituting and summing by parts on $T_2$ using the numerical boundary conditions $(u_0)_{xx} = (u_n)_{xx} = 0$,
    \begin{equation}
        T_2 \;=\; -\frac{D_0}{\Delta x} \sum_i b_{i+1/2}\, D^+ u_i\, \big[(u_{i+1})_{xx} - (u_i)_{xx}\big]\,\Delta x \;=\; D_0 \Delta x \sum_i \big[b_{i+1/2} D^+ u_i - b_{i-1/2} D^+ u_{i-1}\big](u_i)_{xx}.
    \end{equation}
    Splitting the bracket as $b_{i+1/2}(D^+ u_i - D^+ u_{i-1}) + (b_{i+1/2} - b_{i-1/2}) D^+ u_{i-1}$ and using $D^+ u_i - D^+ u_{i-1} = \Delta x\, (u_i)_{xx}$ on a uniform grid,
    \begin{equation}
        T_2 \;=\; \underbrace{D_0 \Delta x^2 \sum_i b_{i+1/2}\, \big((u_i)_{xx}\big)^2}_{T_2^{(a)} \,\geq\, 0} \;+\; \underbrace{D_0 \Delta x \sum_i (b_{i+1/2} - b_{i-1/2})\, D^+ u_{i-1}\, (u_i)_{xx}}_{T_2^{(b)}}. \label{eq:t2-decomp}
    \end{equation}

    \emph{Combined estimate.}
    For $T_2^{(a)}$, substitute the expansion-model $b_{i+1/2} = \Delta x^2 R_{i+1/2}^2/8$ and the inverse inequality $((u_i)_{xx})^2 \leq (2/\Delta x^2)\big((D^+ u_i)^2 + (D^+ u_{i-1})^2\big)$ (from $(u_i)_{xx} = (D^+ u_i - D^+ u_{i-1})/\Delta x$ and $(p-q)^2 \le 2(p^2 + q^2)$):
    \begin{equation}
        T_2^{(a)} \;=\; \frac{D_0 \Delta x^4}{8} \sum_i R_{i+1/2}^2 \big((u_i)_{xx}\big)^2 \;\leq\; \frac{D_0 \Delta x^2}{4} \sum_i R_{i+1/2}^2 \big[(D^+ u_i)^2 + (D^+ u_{i-1})^2\big].
    \end{equation}
    The numerical conditions $(u_0)_{xx} = (u_n)_{xx} = 0$ restrict the cell sum to interior cells $i \in \{1, \ldots, n-1\}$, on which both $D^+ u_i$ (at face $i + 1/2$) and $D^+ u_{i-1}$ (at face $i - 1/2$) are well-defined interior-face forward differences. Bounding $R_{i+1/2}^2 \leq \|R\|_\infty^2$ pointwise and applying the change of summation index $j = i-1$ to the second sum,
    \begin{equation}
        \sum_{i=1}^{n-1} R_{i+1/2}^2 \big[(D^+ u_i)^2 + (D^+ u_{i-1})^2\big] \;\leq\; \|R\|_\infty^2 \bigg[ \sum_{i=1}^{n-1} (D^+ u_i)^2 + \sum_{j=0}^{n-2} (D^+ u_j)^2 \bigg].
    \end{equation}
    Both bracketed sums consist of non-negative terms and are subsets of the interior-face sum $S := \sum_{k=0}^{n-1}(D^+ u_k)^2$, so each is bounded by $S$. Therefore
    \begin{equation}
        T_2^{(a)} \;\leq\; \frac{D_0 \Delta x^2}{2}\, \|R\|_\infty^2\, S. \label{eq:t2a-bound-clean}
    \end{equation}
    The lower bound $|T_1| = D_0 \Delta x \sum_{k=0}^{n-1} a_{k+1/2}(D^+ u_k)^2 \geq D_0 \Delta x\, R_{\min}^2\, S$ (using $a \geq R^2 \geq R_{\min}^2$) gives $S \leq |T_1|/(D_0 \Delta x\, R_{\min}^2)$, and
    \begin{equation}
        T_2^{(a)} \;\leq\; \frac{\Delta x\, \|R\|_\infty^2}{2 R_{\min}^2}\, |T_1|.
    \end{equation}

    For $T_2^{(b)}$, $b_{i+1/2} - b_{i-1/2} = (\Delta x^2/8)(R_{i+1/2}^2 - R_{i-1/2}^2)$, so on a uniform grid the mean value theorem (applied to $R^2$ over $[x_{i-1/2}, x_{i+1/2}]$) gives $|R_{i+1/2}^2 - R_{i-1/2}^2| \leq 2 \|R R_x\|_\infty \Delta x \leq 2\|R\|_\infty \|R_x\|_\infty \Delta x$, hence
    \begin{equation}
        |b_{i+1/2} - b_{i-1/2}| \;\leq\; \frac{\Delta x^3}{4}\, \|R\|_\infty\, \|R_x\|_\infty.
    \end{equation}
    Then
    \begin{equation}
        |T_2^{(b)}| \;\leq\; \frac{D_0 \Delta x^4}{4}\, \|R\|_\infty\, \|R_x\|_\infty \sum_i |D^+ u_{i-1}|\, |(u_i)_{xx}|.
    \end{equation}
    Apply Cauchy--Schwarz to $|T_2^{(b)}|$:
    \begin{equation}
        \sum_i |D^+ u_{i-1}|\, |(u_i)_{xx}| \;\leq\; \bigg(\sum_i (D^+ u_{i-1})^2\bigg)^{1/2} \bigg(\sum_i ((u_i)_{xx})^2\bigg)^{1/2}.
    \end{equation}
    The first factor is bounded by $\sqrt{S}$ (by the change of summation index $j = i-1$, the sum equals $\sum_{j=0}^{n-2}(D^+ u_j)^2 \leq S$). The second factor is bounded via the inverse inequality used for $T_2^{(a)}$,
    \begin{equation}
        \sum_i ((u_i)_{xx})^2 \;\leq\; \frac{2}{\Delta x^2} \bigg[\sum_i (D^+ u_i)^2 + \sum_i (D^+ u_{i-1})^2\bigg] \;\leq\; \frac{4}{\Delta x^2}\, S,
    \end{equation}
    whose square root is $(2/\Delta x)\sqrt{S}$. Combining,
    \begin{equation}
        |T_2^{(b)}| \;\leq\; \frac{D_0 \Delta x^4}{4}\, \|R\|_\infty\, \|R_x\|_\infty \cdot \sqrt{S} \cdot \frac{2\sqrt{S}}{\Delta x} \;=\; \frac{D_0 \Delta x^3}{2}\, \|R\|_\infty\, \|R_x\|_\infty\, S \;\leq\; \frac{\Delta x^2}{2}\, \frac{\|R\|_\infty \|R_x\|_\infty}{R_{\min}^2}\, |T_1|.
    \end{equation}

    \emph{Conclusion.} Combining the two bounds,
    \begin{equation}
        T_2 \;\leq\; T_2^{(a)} + |T_2^{(b)}| \;\leq\; \frac{\Delta x}{2 R_{\min}^2} \left[ \|R\|_\infty^2 + \Delta x\, \|R\|_\infty \|R_x\|_\infty \right] |T_1|.
    \end{equation}
    For the bracketed coefficient on $|T_1|$ to be at most $1$, it suffices that \newline $\Delta x \|R\|_\infty^2/(2 R_{\min}^2) + \Delta x^2 \|R\|_\infty \|R_x\|_\infty/(2 R_{\min}^2) \leq 1$. A sufficient (and slightly stronger) condition is $\Delta x (\|R\|_\infty^2 + \|R\|_\infty \|R_x\|_\infty)/R_{\min}^2 \leq 1$, which is exactly the threshold \eqref{eq:stab-threshold}. Under this condition, $T_1 + T_2 \leq 0$, giving $dE/dt \leq 0$.
\end{proof}

\end{document}
